\documentclass[11pt]{article}

\usepackage[T1]{fontenc}
\usepackage[utf8]{inputenc}
\usepackage{lmodern}
\usepackage{amsmath,amssymb,amsthm,mathtools}
\usepackage{booktabs}
\usepackage{geometry}
\usepackage{microtype}
\usepackage[colorlinks=true,allcolors=blue]{hyperref}

\geometry{margin=1in}

\newtheorem{theorem}{Theorem}[section]
\newtheorem{proposition}[theorem]{Proposition}
\newtheorem{corollary}[theorem]{Corollary}
\newtheorem{lemma}[theorem]{Lemma}
\newtheorem{definition}[theorem]{Definition}
\newtheorem{remark}[theorem]{Remark}

\newcommand{\norm}[1]{\left\lVert #1\right\rVert}
\newcommand{\RR}{\mathbb{R}}
\newcommand{\id}{\operatorname{Id}}
\newcommand{\spec}{\operatorname{spec}}
\newcommand{\tr}{\operatorname{tr}}

\title{Structure-Preserving Reduction of Finite-Dimensional N-Ary Laws:\\
Exact Functoriality, Associator Stability, and Projection Error Bounds}
\author{Eliuth Chavero Jasso\\
\small Independent Researcher\\
\small Apizaco, Tlaxcala, Mexico\\
\small ORCID and email: to be supplied}
\date{Working-paper reconstruction, 29 July 2026}

\begin{document}
\maketitle

\begin{abstract}
We study the reduction of a bounded multilinear $n$-ary law on a finite-dimensional inner-product space to a subspace represented by an isometry. Under exact invariance, the reduction commutes with every planar partial composition, and consequently preserves every multilinear polynomial identity generated by the law. Under approximate closure, we prove an explicit tree-level bound: for a tree with $k$ internal vertices, operator norm $M$ of the law, and closure residual at most $\varepsilon$, the projected composition error is at most $kM^{k-1}\varepsilon$ times the product of the input norms. The associator case gives the concrete constant $2M\varepsilon$. We also state a spectral-snapping perturbation bound under a gap $\gamma$ around the threshold $1/2$, together with a no-gap counterexample. The accompanying SEION Math Core implementation records these results separately from numerical evidence, deduplicates scientific run identities, and fails closed for release while the theorem and multi-seed tightness programs remain incomplete. Thus the paper is a finite-dimensional theorem note and reproducibility specification, not a continuum or universal result.
\end{abstract}

\noindent\textbf{Keywords:} multilinear laws; $n$-ary algebras; operadic composition; associator; invariant subspace; projection error; spectral projector.\\
\noindent\textbf{2020 MSC:} 15A69; 16S37; 47A58; 65F30; 68W30.

\section{Introduction}

Projection is a natural way to reduce a finite-dimensional algebraic object, but it does not automatically preserve the identities that define that object. This distinction is especially important for nonassociative and $n$-ary laws, where the placement of an inner operation is part of the datum. The software release that motivated this paper contains typed ternary conventions, associator diagnostics, projector experiments, spectral controls, and a reproducibility ledger. The present manuscript isolates the theorem-level core and separates it from the empirical claims.

The main question is elementary to state. Let $V$ be a finite-dimensional inner-product space, let $\mu:V^n\to V$ be bounded and multilinear, and let $Q:W\to V$ be an isometry with $P=QQ^*$. If $\mu(QW,\ldots,QW)\subseteq QW$, does the reduced law
\[
    \bar\mu(z_1,\ldots,z_n)=Q^*\mu(Qz_1,\ldots,Qz_n)
\]
retain compositions and identities? The answer is yes, by a direct but useful functoriality argument. If closure is only approximate, the same argument becomes a quantitative telescoping estimate rather than an exact equality.

The contributions are deliberately narrow:
\begin{enumerate}
    \item an exact invariant-reduction theorem for all planar partial compositions;
    \item a polynomial-identity corollary and an explicit associator bound under approximate closure;
    \item a tree-level projection-error bound with its constant computed in terms of $M$, $k$, and $\varepsilon$;
    \item a spectral-snapping stability statement and a sharp discontinuity example at zero gap;
    \item a repository contract that records which statements are proved, tested, observed, conjectural, or blocked.
\end{enumerate}

\section{Relation to prior work}

The $n$-Lie setting and its fundamental identity originate in the classical literature on Filippov algebras \cite{Filippov1985}; broader nonassociative loop and algebraic structures are treated by Sabinin \cite{Sabinin1999}. Operads provide the natural language for composing operations while retaining insertion positions \cite{LodayVallette2012}. Tensor decompositions and identifiability issues motivate the computational representation layer, but they do not by themselves give the reduction theorem proved here \cite{KoldaBader2009,Kruskal1977}.

The perturbation statement uses the standard spectral-subspace viewpoint of Davis--Kahan \cite{DavisKahan1970}. Projection-based model reduction has a substantial systems literature, including structure-preserving methods for dynamical models \cite{BennerGugercinWillcox2015,ChaturantabutBeattieGugercin2016}; the present result is narrower, algebraic, and finite-dimensional. Finite discrete Hodge constructions are supporting examples for the separate cohomology checks, not a novelty claim \cite{Dodziuk1976}.

\begin{table}[t]
\centering
\small
\caption{Conservative prior-art positioning.}
\label{tab:prior-art}
\begin{tabular}{p{0.21\linewidth}p{0.24\linewidth}p{0.24\linewidth}p{0.21\linewidth}}
\toprule
Area & Established result & Exact overlap & Difference and novelty limit \\
\midrule
$n$-ary algebras & Filippov identities and examples & The same multilinear-law setting & No new classification or identity theory is claimed. \\
Operads & Insertion and composition formalism & Planar partial compositions & The contribution is a finite-dimensional reduction lemma, not a new operad. \\
Tensor methods & CP/Tucker/HOSVD representations & Numerical storage and rank sweeps & No new decomposition or identifiability theorem is claimed. \\
Spectral perturbation & Gap-dependent projector estimates & Threshold snapping is a spectral projection & The bound is a direct finite-dimensional specialization. \\
Model reduction & Structure-preserving projection methods & Projection can preserve declared structure & This paper treats an $n$-ary law, not a dynamical-system algorithm. \\
\bottomrule
\end{tabular}
\end{table}

\section{Notation and hypotheses}

All spaces are finite-dimensional real inner-product spaces. The same algebraic arguments hold over $\mathbb C$ with adjoints in place of transposes. The operator norm of $\mu$ is
\[
    M=\norm{\mu}:=\sup_{\norm{x_j}\leq 1}\norm{\mu(x_1,\ldots,x_n)}.
\]
Let $Q:W\to V$ be an isometry, so $Q^*Q=\id_W$, and put $P=QQ^*$. The reduced law is
\[
    \bar\mu=Q^*\mu(Q\cdot,\ldots,Q\cdot):W^n\to W.
\]
For a planar rooted tree $T$ whose internal vertices have arity $n$, write $\mu_T$ for the iterated composition and $k(T)$ for the number of internal vertices. The corresponding reduced tree is $\bar\mu_T$.

\begin{definition}[Closure residual]
For $x=(x_1,\ldots,x_n)\in V^n$, define
\[
    r_\mu(x)=(\id-P)\mu(Px_1,\ldots,Px_n).
\]
The subspace has $(\varepsilon,n)$-approximate closure if
\[
    \norm{r_\mu(x)}\leq \varepsilon\prod_{j=1}^n\norm{x_j}
    \qquad\text{for all }x\in V^n.
\]
\end{definition}

\section{Exact invariant reduction}

\begin{theorem}[Exact functoriality]
\label{thm:exact}
Assume $\mu(QW,\ldots,QW)\subseteq QW$. For every planar $n$-ary tree $T$ and every $z$ in the appropriate product of copies of $W$,
\[
    Q\bar\mu_T(z)=\mu_T(Qz).
\]
In particular, for every partial insertion $\mu\circ_i\mu$,
\[
    Q(\bar\mu\circ_i\bar\mu)(z_1,\ldots,z_{2n-1})
    =(\mu\circ_i\mu)(Qz_1,\ldots,Qz_{2n-1}).
\]
\end{theorem}

\begin{proof}
The statement for one vertex is the definition together with invariance:
\(Q\bar\mu(z)=QQ^*\mu(Qz)=\mu(Qz)\). Suppose it holds for each subtree attached to the root of $T$. Each subtree output belongs to $QW$, so the induction hypothesis writes it as $Q$ times its reduced output. Applying the one-vertex identity at the root gives the claimed equality. The insertion formula is the case $k(T)=2$ with the inner vertex at position $i$.
\end{proof}

\begin{corollary}[Polynomial identities descend]
\label{cor:polynomial}
Let $p=\sum_T c_T T$ be a multilinear polynomial in one $n$-ary operation, with each $T$ a planar tree. If $p(\mu)=0$ on $V$, then $p(\bar\mu)=0$ on $W$.
\end{corollary}

\begin{proof}
Apply Theorem~\ref{thm:exact} term by term. It gives $Qp(\bar\mu)=p(\mu)(Q\cdot)=0$. Since $Q$ is injective, $p(\bar\mu)=0$.
\end{proof}

\section{Approximate closure and explicit constants}

Exact invariance is a qualitative hypothesis. The next theorem quantifies the cost of replacing every intermediate output by its projection. It is stated for arbitrary planar trees so that associators and fundamental-identity residuals are covered by one formula.

\begin{theorem}[Tree-level projection error]
\label{thm:closure}
Assume $\norm{\mu}=M$ and $(\varepsilon,n)$-approximate closure. For a planar $n$-ary tree $T$ with $k=k(T)\geq 1$ internal vertices and leaf inputs $z_1,\ldots,z_{1+k(n-1)}\in W$,
\[
 \norm{Q\bar\mu_T(z)-P\mu_T(Qz)}
 \leq kM^{k-1}\varepsilon\prod_{\ell=1}^{1+k(n-1)}\norm{z_\ell}.
\]
Consequently, for a polynomial $p=\sum_T c_TT$,
\[
 \norm{Qp(\bar\mu)(z)-Pp(\mu)(Qz)}
 \leq \varepsilon\sum_{T:k(T)\geq1}|c_T|k(T)M^{k(T)-1}\prod_\ell\norm{z_\ell}.
\]
\end{theorem}

\begin{proof}
Proceed by induction on the number of internal vertices. At the root, insert and subtract the projected root output. The normal component is bounded by the closure hypothesis. The remaining terms are the differences of the child outputs, each multiplied by a map with norm at most $M$. By the induction hypothesis, every internal vertex contributes one closure residual; a residual below $a$ ancestors is multiplied by at most $M^a$, while the norms of the descendant subtrees contribute the remaining powers of $M$. Thus a tree with $k$ vertices contributes at most $kM^{k-1}\varepsilon$ times the product of leaf norms. Summing the tree estimates with coefficients $c_T$ proves the polynomial statement.
\end{proof}

\begin{corollary}[Associator stability]
\label{cor:assoc}
If an associator is the difference of two ternary trees with two internal vertices and coefficients $+1$ and $-1$, then
\[
 \norm{Q A_{\bar\mu}(z)-P A_\mu(Qz)}
 \leq 2M\varepsilon\prod_{\ell=1}^{5}\norm{z_\ell}.
\]
\end{corollary}

The factor $2M$ is the declared constant for the two-term ternary associator convention. More complicated identities are handled by the polynomial form of Theorem~\ref{thm:closure}; their constants depend on the number and depth of their trees, not on an unspecified generic constant.

\section{Spectral snapping}

Let $A=A^*$ be a symmetric approximate projector and define the threshold map
\[
    \mathcal S(A)=\mathbf 1_{[1/2,\infty)}(A).
\]

\begin{theorem}[Gap stability]
\label{thm:snap}
Suppose $\operatorname{dist}(1/2,\spec(A))\geq\gamma>0$, and let $E=E^*$. If $\norm{E}<\gamma/2$, then
\[
    \norm{\mathcal S(A+E)-\mathcal S(A)}
    \leq \frac{\norm{E}}{\gamma-\norm{E}}
    \leq \frac{2}{\gamma}\norm{E}.
\]
\end{theorem}

\begin{proof}
Weyl's inequality keeps the spectrum of $A+tE$ at distance at least $\gamma-\norm{E}$ from the threshold for $0\leq t\leq1$. The two spectral subspaces separated by $1/2$ therefore remain separated throughout the path. The Davis--Kahan sin-theta estimate applied to this path gives the first inequality; the second follows from $\norm{E}<\gamma/2$.
\end{proof}

\begin{proposition}[No-gap discontinuity]
\label{prop:nogap}
For $A_0=\tfrac12 I_2$ and $E_\delta=\operatorname{diag}(-\delta,\delta)$ with $\delta>0$, the convention above gives $\mathcal S(A_0)=I_2$ and $\mathcal S(A_0+E_\delta)=\operatorname{diag}(0,1)$. Hence
\[
    \norm{\mathcal S(A_0+E_\delta)-\mathcal S(A_0)}=1
    \quad\text{while}\quad \norm{E_\delta}=\delta\to0.
\]
\end{proposition}

\section{Finite cohomology as supporting evidence}

\begin{proposition}[Finite descent]
Let $(C^\bullet,\delta)$ be a finite cochain complex and let $R^k:C^k\to \widetilde C^k$ commute with the differentials: $R^{k+1}\delta^k=\widetilde\delta^kR^k$. Then $R$ maps cycles to cycles, boundaries to boundaries, and induces a linear map $H^k(C)\to H^k(\widetilde C)$. This proposition is a standard supporting fact and is not asserted as the principal novelty of the paper.
\end{proposition}

\begin{proof}
If $\delta^kc=0$, then $\widetilde\delta^kR^kc=R^{k+1}\delta^kc=0$. If $c=\delta^{k-1}b$, then $R^kc=\widetilde\delta^{k-1}R^{k-1}b$. The induced map is therefore well-defined on the quotient of cycles by boundaries.
\end{proof}

\section{Experimental design and reproducibility boundary}

The SEION Math Core run matrix is an evidence harness for finite calculations. The current derived index preserves 73 historical run records and identifies 9 scientific instances after grouping repeated identities; 8 duplicate groups account for 64 repeated records. All 73 records satisfy the declared artifact contract, with 71 complete and 2 runtime failures. These numbers describe software provenance, not independent mathematical experiments.

The current repository also executes fast and full canonical profiles, including a CUDA path on an NVIDIA RTX PRO 5000 Blackwell device. The governance audit records the command, commit, seed, precision, backend, device, artifacts, and limitations. Bound-sharpness experiments with at least five independent seeds, confidence intervals, and tightness ratios remain required before the numerical section can support a release-ready research claim. The distinction follows the project's epistemic policy: a passing verification run is evidence for the registered finite instance, not a proof of a universal theorem.

\section{Limitations and unresolved blockers}

The exact and approximate statements above are finite-dimensional and use a bounded multilinear law. They do not provide a continuum limit, an infinite-dimensional extension, a classification of $n$-ary identities, or a new tensor-identifiability theorem. The spectral result assumes self-adjoint perturbations and a strict threshold gap. The repository's current illustrative curves are not sufficient to estimate asymptotic slopes or bound tightness. ORCID and email metadata are intentionally left pending rather than fabricated.

The machine-readable blockers are B-0001 (the central theorem program must be integrated with the existing implementation), B-0002 (the historical index contains repeated executions and must not be treated as independent evidence), B-0003 (the quantitative figure suite needs registered multi-seed data), and B-0004 (author contact metadata is missing). Release is therefore fail-closed until the audit is green.

\section{Conclusion}

Exact invariant reduction is functorial for planar compositions, and approximate closure produces a computable error budget. The tree-level form makes the dependence on composition depth explicit and specializes to the $2M\varepsilon$ associator bound. Spectral snapping is stable only while a threshold gap is present. These are the theorem-level foundations for the SEION verification suite; the companion manuscript documents the software and evidence contracts separately.

\bibliographystyle{plain}
\bibliography{references}

\end{document}
